% Videos: modular arithmetic parts 1-2; college Lecture 4 (residues, inverses, Diophantine)
\section{Modular Arithmetic}\label{sec:modular}

\begin{keybox}[title=Definitions]
\vspace{-6pt}
\[
a \equiv b \pmod m \iff m \mid (a - b) \iff a = km + b \text{ for some } k \in \mathbb{Z}
\]
The operator $a \bmod m$ always returns the \textbf{least non-negative residue} $0 \le r < m$. \textbf{Residue class} $[a]_m$: all integers $\equiv a \pmod m$. $\mathbb{Z}_n = \{0, \dots, n-1\}$; $\mathbb{Z}_n^* = \{a \in \mathbb{Z}_n : \gcd(a,n) = 1\}$; for a prime $p$, $\mathbb{Z}_p^* = \{1, \dots, p-1\}$.
\end{keybox}

\textbf{Negative numbers}: $-255 = (-23)(11) + (-2)$ has $r < 0$; fix it by $q \to q - 1$, $r \to r + n$: $-255 = (-24)(11) + 9$, so $-255 \bmod 11 = 9$.

\begin{center}\small
\begin{tabular}{@{}llll@{}}
\toprule
Statement & $a - b$ & Divisible by $m$? & \\
\midrule
$38 \equiv 14 \pmod{12}$ & 24 & yes & \textbf{valid} (but $38 \bmod 12 = 2$, not 14)\\
$10 \equiv 2 \pmod 6$ & 8 & no & \textbf{invalid}\\
$13 \equiv 3 \pmod{13}$ & 10 & no & \textbf{invalid}\\
$-8 \equiv 7 \pmod 5$ & $-15$ & yes & \textbf{valid}\\
$-3 \equiv -8 \pmod 5$ & 5 & yes & \textbf{valid}\\
\bottomrule
\end{tabular}
\end{center}

\subsection{Reduce first}
\begin{keybox}
\vspace{-6pt}
\[
(a \pm b) \bmod n = [(a \bmod n) \pm (b \bmod n)] \bmod n \qquad (a b) \bmod n = [(a \bmod n)(b \bmod n)] \bmod n
\]
There is \textbf{no division rule}: dividing by $a$ means multiplying by $a^{-1}$, which exists only if $\gcd(a, n) = 1$.
\end{keybox}
E.g.\ $(123 \times 456) \bmod 7 = (4 \times 1) \bmod 7 = 4$; \; $3^{100} \bmod 10 = (3^4)^{25} \bmod 10 = 1$.

\subsection{Properties of $\mathbb{Z}_n$ and inverses}
$+$ and $\times$ are commutative, associative, distributive; identities 0 and 1. Subtraction is not commutative.
\begin{itemize}
  \item \textbf{Additive inverse}: $b = n - a$ (every element has one). Pairs in $\mathbb{Z}_{10}$: (0,0), (1,9), (2,8), (3,7), (4,6), (5,5).
  \item \textbf{Multiplicative inverse}: $a \cdot a^{-1} \equiv 1 \pmod n$, \textbf{exists iff $\gcd(a, n) = 1$}. Found with the extended Euclidean algorithm (Section~\ref{sec:exteuclid}). $8^{-1}$ in $\mathbb{Z}_{10}$: none. $11^{-1} \bmod 26 = 19$; $23^{-1} \bmod 100 = 87$; $12^{-1} \bmod 26$: none ($\gcd = 2$).
\end{itemize}
\begin{center}\small\setlength{\tabcolsep}{6pt}
\begin{tabular}{@{}l*{8}{c}@{}}
\toprule
$a$ in $\mathbb{Z}_8$ & 0&1&2&3&4&5&6&7\\
$-a$ & 0&7&6&5&4&3&2&1\\
$a^{-1}$ & --&1&--&3&--&5&--&7\\
\bottomrule
\end{tabular}
\end{center}
Caesar decryption $=$ adding the additive inverse of the key; Hill decryption needs $(\det K)^{-1} \bmod 26$, so $\gcd(\det K, 26) = 1$.

\subsection{Linear congruences $ax \equiv b \pmod n$}
Let $d = \gcd(a, n)$. If $d \nmid b$: \textbf{no solution}. If $d \mid b$: \textbf{exactly $d$ solutions}. Divide through by $d$, solve $x_0 = (b/d)(a/d)^{-1} \bmod (n/d)$; the solutions are $x_0 + k(n/d)$, $k = 0, \dots, d-1$.
\begin{itemize}
  \item $10x \equiv 2 \pmod{15}$: $d = 5 \nmid 2$ $\Rightarrow$ no solution.
  \item $14x \equiv 12 \pmod{18}$: $d = 2$; $7x \equiv 6 \pmod 9$, $7^{-1} = 4$, $x_0 = 24 \bmod 9 = 6$; solutions \textbf{6, 15}.
  \item $3x + 4 \equiv 6 \pmod{13}$: $3x \equiv 2$, $3^{-1} = 9$, $x = 18 \bmod 13 = \mathbf{5}$.
\end{itemize}

\begin{examtip}
\begin{itemize}
  \item $a \equiv b \pmod m \iff m \mid (a - b)$: fastest validity test.
  \item ``$\equiv \pmod m$'' allows any representative; ``$\bmod$'' gives only $0..m-1$.
  \item $a^{-1} \bmod n$ exists only when $\gcd(a, n) = 1$; a linear congruence has $\gcd(a,n)$ solutions or none.
\end{itemize}
\end{examtip}
